\section{相関係数とRによる実践}
\label{b06_Corr}

\subsection{授業内容}

\subsubsection{概要}
これまでは1変数の記述統計について学んできたが，本コマからは複数の変数間の関係性について考察する。変数同士の関係性を数値化して表す指標として，共分散や相関係数について学ぶ。また，変数の尺度水準に応じて適切な関連指標が異なることを理解し，それぞれの特徴と適用条件について理解する。相関係数は直線的関係についての指数であることに注意し，散布図による可視化の重要性を理解する。実際のデータを用いてRによる計算と可視化の演習を行うことで，相関分析の実践的技能を身につける。

\subsubsection{コマ主題細目}
\begin{description}
	\item[二変数間関係の指標] 変数同士の関係を記述する方法として，クロス集計表，共分散，相関係数などがある。クロス集計表はカテゴリカルデータの関連を表現する基本的な方法である。共分散は二つの変数がどのように連動して変化するかを数値化したものだが，変数の単位に依存するという欠点がある。そこで単位の影響を取り除くために標準化したものが相関係数であり，-1から+1の範囲で二変数間の線形関係の強さを表す。これらの基本的な考え方と計算方法について理解する。
	      \begin{flushright}
		      $\to$ \textcite{Yamauchi201001}のP.62--68,\textcite{Yamada2004}のP.44--50，\textcite{Kawabata201408}のP.48--52.
	      \end{flushright}

	\item[尺度水準に応じた相関係数] 変数の尺度水準によって適切な関連指標は異なる。名義尺度の関連はカイ2乗値やファイ係数，クラメールのV係数などで表される。二値データ間の関連を表すテトラコリック相関，順序尺度間の関連を表すポリコリック相関，順序尺度と連続変数の関連を表すポリシリアル相関などがある。連続変数間の関連を表す指標としては，ピアソンの積率相関係数が最も一般的であるが，線形関係以外も検出できるスピアマンの順位相関係数やケンドールのタウ，さらに最近では非線形関係も捉えられるMIC(Maximum Information Coefficient)なども開発されている。これらの指標の特徴と適用条件を理解し，データの性質に応じた適切な指標を選択できるようになる。
	      \begin{flushright}
		      $\to$ 様々な相関係数について\textcite{Kawabata201408}のP.53--58，\textcite{Yamada2004}のP.51--59，MICやテトラコリック相関などカテゴリカルな相関係数については\textcite{TDE2022}の2章を参照。
	      \end{flushright}

	\item[Rによる相関分析の実践] 相関係数の計算と検定はRを用いて簡単に行うことができる。\texttt{cor()}関数を用いた相関係数の計算方法について学ぶ。また，\texttt{polycor}パッケージを用いたポリコリック相関やポリシリアル相関の計算方法，\texttt{minerva}パッケージを用いたMICの計算方法についても触れる。実データを用いた演習を通じて，相関分析の実践的な技術を身につける。
	      \begin{flushright}
		      $\to$ Rによる相関分析については\textcite{kosugi2019}のP.110--120を参照。
	      \end{flushright}

	\item[相関の落とし穴と可視化の重要性] 相関係数だけを見て判断することの危険性について理解する。同じ相関係数の値でも，散布図のパターンは大きく異なる可能性がある。この点を示すための有名な例として，アンスコムの例や，より現代的な\texttt{datasauRus}パッケージに含まれるDinosaurus datasetがある。これらの例を通じて，相関分析を行う際には必ず散布図を描いてデータの分布を確認することの重要性を学ぶ。また，外れ値の影響や非線形関係，層別の問題など，相関分析における様々な落とし穴について理解し，それらを回避するための方法について学ぶ。
	      \begin{flushright}
		      $\to$ \verb|datasauRus|パッケージには，記述統計が同じだがプロットがまったく異なる13のケースが含まれていて参考になる。可視化の重要性については\textcite{Kieran2018}を参照。
	      \end{flushright}
\end{description}

\subsubsection{キーワード}
\begin{itemize}
	\item 共分散と相関係数
	\item カイ2乗値，ファイ係数，テトラコリック相関，ポリコリック相関
	\item ピアソンの積率相関係数，スピアマンの順位相関，MIC
	\item 散布図による可視化
\end{itemize}

\subsection{授業情報}
\paragraph{コマの展開方法}
講義と演習
\paragraph{標準シラバスにおける位置づけ}
科目番号5 心理学統計法；1.心理学で用いられる統計手法；B 記述統計：相関
\subsubsection{予習・復習課題}
\paragraph{予習}
分散，標準偏差，標準化の手続きについて再確認しておくこと。分散は自分自身との共分散という形で理解することができるし，共分散を標準化したものが相関係数になるからである。また，Rの基本操作とデータ読み込みの方法について復習しておくこと。
\paragraph{復習}
授業で扱ったサンプルデータを用いて，以下の課題を行い，Rmarkdownファイルにまとめて提出すること。
1. 複数の変数間の相関係数を計算し，相関行列を作成する
2. 相関係数のヒートマップを作成する
3. 散布図を描き，相関係数の値と視覚的な関係性を比較考察する
4. \texttt{datasauRus}パッケージの例を再現し，同じ相関係数でも異なるパターンが生じることを確認する
5. 身近な例から二つの変数を選び，それらの関係を相関分析によって検討する

また，相関関係と因果関係の違いについて理解を深めるため，日常生活や報道などで見られる「相関関係を因果関係と誤解している例」を一つ挙げ，なぜそれが誤りであるかを説明すること。これはデータリテラシーの問題でもあり，統計的思考の基本でもある。